% Chapter 3: Encoder Bias Modeling and POMDP Formulation

\chapter{Encoder Bias Modeling and POMDP Formulation}\label{ch:problem_formulation}

\section{Encoder Calibration Bias in Industrial Robots}\label{sec:bias_sources}

% 5 sources of encoder bias:
% - Factory calibration inaccuracy (fixed, 0.01-0.1 rad)
% - Temperature drift (slow, 0.05-0.2 rad)
% - Uncalibrated motor/reducer replacement (fixed, 0.1-0.5 rad)
% - Post-collision misalignment (sudden, 0.1-1.0 rad)
% - Electromagnetic interference (noise, 0.001-0.01 rad)
%
% This work focuses on fixed/random biases (covers sources 1,3,4)

\section{Causal Chain: One Fault Source, Three Effects}\label{sec:causal_chain}

% Mathematical derivation of three cascading effects:
%
% Given bias vector b \in R^7, the encoder reports:
%   q_measured = q_true + b
%
% Effect 1: Initial position offset
%   PD controller steady state: q_measured = q_home -> q_true = q_home - b
%   The robot thinks it's at home but physically offset
%
% Effect 2: Control (execution) bias
%   OSC computes: tau = J(q_measured)^T * Lambda * F_imp
%   Since J(q_true + b) != J(q_true), torque direction and magnitude are wrong
%   -> end-effector moves in wrong direction
%
% Effect 3: Perception bias
%   FK reports: x_hat_ee = FK(q_measured) = FK(q_true + b)
%   Since FK(q+b) != FK(q), the reported TCP position is wrong
%   -> policy receives incorrect state information
%
% Key: all three effects stem from the SAME bias vector b

\section{POMDP Formalization}\label{sec:pomdp}

% 3.3.1 Standard MDP (no bias)
%   M = (S, A, T, R, gamma)
%   S = (q, dq, g, x_ee, x_obj)  -- fully observable
%   Policy: pi(a|s)
%
% 3.3.2 POMDP under encoder bias
%   M = (S, A, T, R, Omega, O, gamma)
%   S = (q_true, dq, g, x_ee_true, x_obj, b)  -- b is latent
%   Omega = observation space
%   O(o|s): observation function
%     o = (q_true + b, dq, g, FK(q_true + b), x_obj)
%
%   Key properties of b:
%   - Episode-level constant: b' = b (does not change within episode)
%   - Sampled per episode: b ~ P(b)
%   - Not directly observable by the agent
%
% 3.3.3 Distinction from standard POMDP
%   - Hidden variable b is static within episode (not evolving)
%   - Transition function: T(s'|s,a) preserves b
%   - This is closer to a contextual MDP / hidden-parameter MDP

\section{Observability Analysis}\label{sec:observability}

% 3.4.1 Unidentifiability under standard observations (18D)
%
%   Proposition: Given single-step observation o = (q+b, dq, g, FK(q+b), x_obj),
%   there exist infinitely many (q, b) pairs producing the same o.
%   -> The POMDP is not identifiable from single-step proprioception alone.
%
%   Proof sketch: o_joint = q + b. For any delta, (q-delta, b+delta) produces
%   the same o_joint. FK(q+b) depends only on the sum, not individual q,b.
%
%   This explains the 18D ablation failure (42% degradation).
%
% 3.4.2 Observation augmentation with external position signal
%
%   Augmented observation: o' = o \cup {x_ee_true + epsilon}
%   where epsilon ~ N(0, sigma^2 I) is sensor noise (sigma ~ 5mm)
%
%   Key signal: Delta_x = real_tcp - biased_tcp
%             = FK(q_true) - FK(q_true + b) + epsilon
%             ~ f(b) + epsilon
%
%   Delta_x is an implicit observation of b
%   -> With sufficiently small epsilon, (o, real_tcp) can approximately
%      determine both q_true and b
%   -> POMDP approximately reduces to MDP
%
% 3.4.3 Role of noise magnitude
%   - sigma -> 0: perfect identifiability, policy can fully compensate
%   - sigma -> inf: real_tcp becomes uninformative, back to POMDP
%   - sigma = 5mm: practical regime, sufficient for bias identification
%
% 3.4.4 Why not just use real_tcp alone?
%   Biased proprioception still carries information about joint configuration
%   (e.g., singularity proximity, joint limits) that real_tcp alone cannot provide.
%   The policy needs BOTH signals.

\section{Comparison with Alternative POMDP Solutions}\label{sec:pomdp_comparison}

% Belief state methods: maintain P(b|o_{1:t}) -> computationally expensive
% RNN/GRU implicit inference: needs multi-step history -> lower sample efficiency
% Explicit bias estimation (future work): train a bias estimator network
% Our approach: observation augmentation -> single-step identifiability
%   -> can directly use MDP algorithms (SAC)
%   Advantages: simple, no architectural changes, sample efficient

\section{Simulation Implementation}\label{sec:sim_implementation}

% How the three effects are implemented in MuJoCo:
%
% Effect 1 (init offset): q_init = q_home - b during reset
% Effect 2 (control bias): temporarily replace qpos with q+b before OSC,
%   compute tau, then restore true qpos before mj_step
% Effect 3 (perception bias): get_robot_state() returns q+b and FK(q+b)
%
% External sensor simulation:
% - Block position: from MuJoCo sensor (external camera, unaffected by bias)
% - Noisy real TCP: true FK + N(0, 5mm) noise (simulates ArUco/OptiTrack)
%
% Simulation-to-reality gap analysis (8 items)
% -> simulation is conservative: bias effects are smaller than reality
